\chapter{Implementacija}
\label{ch:implementacija}

V tem poglavju je predstavljena izvedba celotne poti od zvoka v prostoru do
odločitve, prikazane na zaslonu razvojne plošče. Končni program deluje brez
operacijskega sistema in brez omrežne povezave. Procesorsko jedro zato poleg
zagona aplikacije upravlja zajem zvočnih vzorcev, predobdelavo, odločanje,
zaslon in diagnostični izhod, izračun nevronske mreže pa je v največji možni
meri prenesen na nevronski pospeševalnik Neural-ART.

Kot preverjeno strojno izhodišče smo uporabili uradni primer za obdelavo zvoka
na družini STM32N6 \cite{ststm32n6audio}. Primer je že vseboval gonilnike za
vgrajeni mikrofon, knjižnico za izračun zvočnih značilk in povezavo z izvajalnim
okoljem ST. Nad to osnovo smo zamenjali model, prilagodili parametre
predobdelave, dodali šestrazredno odločanje, časovno filtriranje, lasten
uporabniški vmesnik in strukturirano beleženje meritev. Tak pristop je zmanjšal
tveganje, da bi napako v osnovni nastavitvi strojne opreme zamenjali za napako
modela.

\section{Zajem in prenos zvočnih podatkov}

Uporabljena razvojna plošča ima vgrajen digitalni mikrofon IMP34DT05 in
petpalčni zaslon, zato za delovanje prototipa ni potrebna dodatna zvočna
strojna oprema \cite{st2024stm32n6570dk}. Mikrofon oddaja pulzno-gostotni
digitalni signal. Tak signal še ni zaporedje običajnih šestnajstbitnih zvočnih
vzorcev, ki jih pričakuje predobdelava, zato ga je treba najprej filtrirati in
zmanjšati njegovo vzorčno hitrost.

V končni izvedbi smo za ta korak uporabili prvi večnamenski digitalni filter
mikrokrmilnika. Gonilnik plošče nastavi vir na digitalni mikrofon, vzorčno
frekvenco na \SI{16000}{\hertz}, ločljivost na 16 bitov in število kanalov na
ena. Večnamenski digitalni filter opravi strojno filtriranje ter iz
pulzno-gostotnega toka pripravi 32-bitne vmesne vrednosti. Gonilnik vsako
vrednost deli z 256, jo omeji na dovoljeno območje predznačenega
šestnajstbitnega števila in jo shrani kot enokanalni zvočni vzorec. Na tej
točki je signal po obliki primerljiv z običajnim enokanalnim zvočnim zapisom.

Prenos iz periferne enote v pomnilnik poteka z neposrednim dostopom do
pomnilnika. Uporabljena je krožna prenosna vrsta, zato procesorskemu jedru ni
treba prebrati vsakega vzorca posebej. Zajemni medpomnilnik vsebuje 320
vzorcev, kar pri \SI{16000}{\hertz} pomeni \SI{20}{\milli\second} zvoka.
Prenosna enota sproži povratni klic po zapolnitvi prve in druge polovice
medpomnilnika. Posamezen povratni klic tako preda 160 vzorcev oziroma
\SI{10}{\milli\second} zvoka.

Povratni klici morajo biti kratki, saj se izvajajo med neprekinjenim zajemom.
Njihova naloga je le kopiranje nove polovice medpomnilnika v večji krožni
medpomnilnik aplikacije. Ta hrani položaj naslednjega zapisa, položaj
naslednjega branja in število razpoložljivih vzorcev. Pri doseženem koncu
pomnilniškega območja se položaj vrne na začetek. Med spremembo skupnih
kazalcev so prekinitve za kratek čas onemogočene, s čimer se prepreči, da bi
povratni klic spremenil stanje med branjem glavne zanke.

Glavna zanka čaka, da se v krožnem medpomnilniku zbere 15.360 novih vzorcev.
To ustreza \SI{960}{\milli\second} novega zvoka. Pred njih se ohrani še 240
vzorcev prejšnjega bloka, ker se zadnje časovno okno prejšnje značilke in prvo
časovno okno nove značilke prekrivata. Predobdelava zato prejme 15.600 vzorcev,
ki obsegajo \SI{975}{\milli\second} zvočnega signala. Takšna dolžina ni
naključna: prvo časovno okno vsebuje 400 vzorcev, preostalih 95 oken pa se
začne v razmiku po 160 vzorcev,

\begin{equation}
  400 + 95 \cdot 160 = 15\,600.
  \label{eq:dolzina-zvocnega-bloka}
\end{equation}

Zajem in obdelava sta s tem razvezana. Periferna enota ves čas sprejema zvok,
glavna zanka pa obdeluje le dovolj velike bloke. Če bi predobdelava neposredno
brala iz majhnega zajemnega medpomnilnika, bi ga lahko strojni prenos prepisal,
še preden bi bila obdelava končana. Krožni medpomnilnik prepreči takšno
tekmovanje za podatke in omogoča neprekinjeno delovanje brez operacijskega
sistema.

\section{Predobdelava zvočnega signala}

Nevronska mreža ne prejme surovih zvočnih vzorcev. Njena vhodna predstavitev
je logaritemski melov spektrogram, ki prikaže spreminjanje frekvenčne vsebine
skozi čas. Parametri morajo biti enaki parametrom pri učenju modela. Že majhna
razlika, na primer druga vzorčna frekvenca ali drugačna razporeditev melovih
filtrov, bi povzročila drugačno vhodno porazdelitev in slabše napovedi.

Za vsak vhodni blok izvedemo naslednje korake. Signal razdelimo na 96 oken po
400 vzorcev. Dolžina posameznega okna je \SI{25}{\milli\second}, začetek
naslednjega okna pa je zamaknjen za 160 vzorcev oziroma
\SI{10}{\milli\second}. Vsako okno pomnožimo s Hannovo okensko funkcijo in ga
simetrično dopolnimo z ničlami na 512 vzorcev; na vsaki strani dodamo 56 ničel.
Nato z realno Fourierjevo transformacijo izračunamo magnitudni spekter.

Frekvenčne koše preslikamo v 64 melovih pasov med \SI{125}{\hertz} in
\SI{7500}{\hertz}. Uporabljena je melova preslikava po formuli, ki jo uporablja
orodje Hidden Markov Model Toolkit. Nad energijami melovih pasov uporabimo
naravni logaritem. Rezultat enega zvočnega bloka je matrika s 64 frekvenčnimi
pasovi in 96 časovnimi stolpci. To je tudi vhodna predstavitev, predpisana za
družino modelov YAMNet \cite{tensorflowyamnetrepo}.

Tabela \ref{tab:predobdelava-parametri} povzema parametre dejanske programske
konfiguracije. Vrednosti niso le opis predvidene zasnove, temveč so prepisane
iz konfiguracije končne izdaje programa.

\begin{table}[htbp]
  \centering
  \caption{Parametri predobdelave zvočnega signala.}
  \label{tab:predobdelava-parametri}
  \renewcommand{\arraystretch}{1.12}
  \begin{tabular}{>{\raggedright\arraybackslash}p{0.47\textwidth}>{\raggedright\arraybackslash}p{0.36\textwidth}}
    \toprule
    Parameter & Vrednost \\
    \midrule
    Vzorčna frekvenca & \SI{16000}{\hertz}, en kanal \\
    Dolžina časovnega okna & 400 vzorcev oziroma \SI{25}{\milli\second} \\
    Razmik med okni & 160 vzorcev oziroma \SI{10}{\milli\second} \\
    Dolžina Fourierjeve transformacije & 512 vzorcev \\
    Število časovnih stolpcev & 96 \\
    Število melovih pasov & 64 \\
    Frekvenčno območje & od \SI{125}{\hertz} do \SI{7500}{\hertz} \\
    Vrsta spektra & magnitudni spekter \\
    Frekvenčna preslikava & melova preslikava po formuli orodja Hidden Markov Model Toolkit \\
    Logaritemska preslikava & naravni logaritem \\
    Oblika vhodnega tenzorja & en primer, en kanal, 96 časovnih korakov, 64 melovih pasov \\
    Vrsta vhodnih elementov & predznačena osembitna cela števila \\
    \bottomrule
  \end{tabular}
\end{table}

Matrika se že med izračunom pretvori v obliko, ki jo pričakuje model. Vsak
realni element $x$ delimo s kvantizacijskim korakom, rezultat zaokrožimo na
najbližje celo število, mu prištejemo ničelno točko in ga omejimo na območje
predznačenega osembitnega števila:

\begin{equation}
  q = \max\!\left(-128,\min\!\left(127,
      \left\lfloor\frac{x}{0{,}034300160}\right\rceil + 51\right)\right).
  \label{eq:kvantizacija-vhoda}
\end{equation}

Število $0{,}034300160$ je kvantizacijski korak, število 51 pa ničelna točka
vhodnega tenzorja. Obe vrednosti sta pridobljeni neposredno iz prevedenega
modela ob inicializaciji, zato v kodi predobdelave nista podvojeni kot ročno
vpisani konstanti. Končni model uporablja vrstni red
\texttt{[1, 1, 96, 64]}. Vseh 64 melovih vrednosti enega časovnega koraka je
zato v pomnilniku zapisanih zaporedno, nato sledijo vrednosti naslednjega
časovnega koraka. Ta razporeditev se razlikuje od zgodnejše različice modela in
jo je bilo treba uskladiti tudi v vgrajeni predobdelavi. Po zapisu se območje
podatkovnega predpomnilnika očisti in razveljavi. Tako nevronski pospeševalnik
vidi najnovejše podatke in ne stare kopije, ki bi še ostala v predpomnilniku
procesorskega jedra.

Predobdelava hkrati izračuna vsoto vrednosti spektrograma. Če vsota ne preseže
praga 4000, se blok obravnava kot neaktiven. Model se sicer izvede tudi za tak
blok, vendar odločitveni filter rezultata ne sme potrditi kot dogodek. Prag s
tem deluje kot preprost vhodni nadzor aktivnosti in preprečuje, da bi tišina
neprestano ustvarjala razred z največjo, vendar vseeno nezanesljivo oceno.

\section{Priprava in pretvorba modela}

Izvirni YAMNet uporablja arhitekturo MobileNetV1 in vrača ocene za 521
zvočnih dogodkov iz zbirke AudioSet \cite{tensorflowyamnet,hershey2017cnn}.
Tak model je uporaben kot splošni razpoznavalnik zvoka, vendar njegov izhod ni
prilagojen izbranemu naboru nevarnosti. V končni izvedbi smo zato uporabili
predhodno naučeno različico YAMNet-1024 kot izločevalnik značilk. Število 1024
označuje dolžino vektorja značilk pred zaključno klasifikacijsko plastjo. Ta
različica je večja od prvotno preskušene različice YAMNet-256, vendar jo lahko
namenski nevronski pospeševalnik na mikrokrmilniku STM32N6 še vedno izvaja v
realnem času.

Končni model ima šest izhodov: pasji lajež, razbitje stekla, strel oziroma streljanje,
sireno, govor in nevihto. Pet razredov predstavlja dogodke, ki jih sistem
obravnava kot možno nevarnost, govor pa je informacijski oziroma zaščitni razred.
Dodali smo ga zato, ker se je v zgodnjih preskusih pokazalo, da glasen človeški
glas lahko po zvočni strukturi tekmuje z nekaterimi nevarnostnimi razredi.
Ločen izhod za govor omogoča, da sistem tak vhod poimenuje in ga ne spremeni v
nevarnostni alarm.

Učenje je potekalo na osebnem računalniku, ne na mikrokrmilniku. Izhodiščni
uravnoteženi učni nabor je vseboval po 192 posnetkov za vsak razred, skupaj
1.152 posnetkov iz zbirk ESC-50 in FSD50K. Pri razbitju stekla smo široko
oznako \emph{Glass} nadomestili z izrecno oznako \emph{Shatter}. S tem smo
izločili primere rokovanja, cingljanja in udarcev ob steklo, pri katerih do
razbitja sploh ne pride. Izbira posnetkov je upoštevala tudi izvornega
nalagatelja, zato se viri za učenje in razvojno preverjanje pri tem razredu ne
prekrivajo.

Kratka dogodka, kakršna sta strel in razbitje stekla, bi se v daljšem posnetku
lahko izgubila med številnimi spektrogramskimi izseki brez ciljnega dogodka.
Zato smo zanju izdelali posebno povečanje učnih podatkov. Program najprej
poišče čas dogodka, nato izreže enosekundno okno in ga z majhnim naključnim
zamikom postavi na različna mesta. Pri strelih uporabi najmočnejše
20-milisekundno energijsko območje ter ohrani najmanj 620 milisekund signala po
dogodku. Če se strel pojavi na samem začetku posnetka, manjkajoči del okna
dopolni z ničlami na levi strani in ne zavrže začetka dogodka. Pri steklu se
uporabi iskalnik značilnega energijskega in spektralnega prehoda, dodatno
varovalo pa sumljivo tih kandidat zamenja z najmočnejšim kratkotrajnim vrhom.

Iz vsakega izbranega prehodnega izvora smo pripravili šest različic. Spremembe
vključujejo položaj dogodka, zmerno spremembo ojačenja, blago frekvenčno
filtriranje, ponavljanje in dodajanje tihega ozadja iz varnih razredov zbirke
ESC-50. Za razbitje stekla nismo spreminjali hitrosti in nismo uporabili močnega
stiskanja amplitude, ker sta takšni spremembi pri prejšnjem poskusu kvarili
naravni zven dogodka. Zaščita vrha po potrebi oslabi celotno mešanico in je ne
trdo odreže. Deterministično izbranih 24 povečanih primerov strelov in razbitja
stekla smo pred učenjem poslušali. Vsi so vsebovali celovit dogodek, 23 pa jih
je bilo ocenjenih kot povsem značilnih. Šele po uspešnem pregledu je skripta
dovolila začetek učenja.

Končni učni manifest vsebuje 4.032 vrstic: 384 za pasji lajež, po 1.344 za
razbitje stekla in strel, 192 za sireno, 576 za govor ter 192 za nevihto.
Razvojni preveritveni nabor je ostal nespremenjen in vsebuje 278 posnetkov:
po 60 posnetkov pasjega laježa, razbitja stekla in nevihte, 54 posnetkov
strelov, 20 posnetkov siren ter 24 posnetkov govora. Povečane različice so bile
uporabljene samo za učenje in so podedovale skupino svojega izvornega posnetka,
zato različica istega dogodka ni mogla preiti v razvojni preveritveni del.

Med učenjem so bili zvočni posnetki pretvorjeni v enake logaritemske melove
spektrograme kot na napravi. Predhodno naučeni del YAMNet-1024 je ostal
nespremenjen, učila pa se je nova polno povezana zaključna plast s šestimi
izhodi. Celoten model ima 3.223.494 parametrov, od tega je 6.150 učljivih.
Uporabljen je bil optimizator Adam z začetno hitrostjo učenja 0,001, velikostjo
učne skupine 16 in največ 60 prehodi čez učne podatke. Ob nespremenjeni
razvojni pravilnosti se je hitrost učenja zmanjševala, zgodnja ustavitev pa je
obnovila najboljše uteži.

Po učenju smo izvedli kvantizacijo po učenju. Uteži in vmesni tenzorji
podprtih operacij so zapisani z osembitnimi celimi števili, zunanji vhod modela
je predznačen osembitni tenzor, šest izhodnih ocen pa ostane v 32-bitnem zapisu
s plavajočo vejico. Končni model je shranjen v obliki Open Neural Network Exchange z
izrecnimi operacijami kvantizacije in dekvantizacije. Datoteka je velika
3.462.018 bajtov in je označena z 256-bitno kriptografsko kontrolno vsoto, kar omogoča
nedvoumno povezavo med učenjem, pretvorbo in različico na napravi.

Za namestitev smo uporabili ST Edge AI Core, uradno orodje za analizo,
pretvorbo in izvajanje modelov na napravah STM32 \cite{st2026edgeai}. Različica
4.0.1 je iz modela pripravila statične opise omrežja in programsko kodo v
jeziku C, razporedila operacije med procesorsko jedro in Neural-ART ter
ustvarila ločeno binarno datoteko uteži velikosti 3.279.505 bajtov. Ločitev
aplikacije in uteži omogoča njuno zapisovanje v različni področji zunanjega
bliskovnega pomnilnika.

Poročilo prevajalnika navaja 72.748.547 operacij množenja in seštevanja za en
prehod modela. Ta podatek opisuje računsko zahtevnost omrežja, ne pa časa
izvajanja. Dejanski čas je odvisen od razporeditve operacij, dostopov do
pomnilnika, frekvence in uporabljenega pospeševalnika, zato je izmerjen ločeno
v poglavju o rezultatih.

\section{Izvajanje modela na Neural-ART}

Mikrokrmilnik STM32N6 združuje procesorsko jedro Arm Cortex-M55 in namenski
nevronski pospeševalnik Neural-ART. Programska orodja ST so namenjena prav
pretvorbi modelov in pripravi razporeda za ta pospeševalnik
\cite{st2026stm32n6ai}. V naši izvedbi procesorsko jedro še vedno vodi celoten
program, Neural-ART pa izvede računsko najzahtevnejše plasti nevronske mreže.

Ob zagonu aplikacija inicializira izvajalno okolje, ustvari kontekst omrežja in
preveri število ter obliko vhodnega in izhodnega tenzorja. Nato pridobi kazalca
na vhodni in izhodni medpomnilnik ter iz opisa vhodnega tenzorja prebere
kvantizacijski korak in ničelno točko. Predobdelava zato piše neposredno v
vhodno območje, ki ga uporablja prevedeno omrežje; dodatno kopiranje celotnega
spektrograma ni potrebno.

Ko je vhod pripravljen, aplikacija pokliče sinhrono izvajanje omrežja. Z vidika
glavne zanke to pomeni, da funkcija vrne nadzor šele, ko so na voljo vsi šesti
izhodi. Tak način poenostavi nadaljnje odločanje, ker odločitveni filter ne more
prebrati delno izračunanega rezultata. Čas izvajanja se meri s števcem ciklov
procesorskega jedra in se shrani med diagnostične merilne podatke.

Prevedeni razpored je razdeljen na 34 zaporednih odsekov izvajanja. Devetindvajset
odsekov se izvede v celoti strojno, pet pa programsko. Iz števila odsekov ni
mogoče izračunati odstotka strojnega pospeševanja, saj posamezni odseki ne
vsebujejo enakega števila operacij. Podatek vseeno potrjuje, da večino
strukture omrežja prevajalnik lahko preslika na Neural-ART, hkrati pa ohrani
programsko pot za nepodprte ali pomožne operacije.

Za začasne rezultate omrežja je prevajalnik rezerviral 245.760 bajtov
notranjega pomnilnika, kar vključuje tudi vhodni in izhodni medpomnilnik.
Programska koda, zajem zvoka, predobdelava, filtriranje odločitev in risanje
zaslona se izvajajo na procesorskem jedru. Takšna delitev nalog je bistvo
pospešene izvedbe: pospeševalnik ne nadomesti celotnega programa, ampak prevzame
plastne operacije nevronske mreže, medtem ko procesorsko jedro skrbi za vhod,
nadzor in uporabniško izkušnjo.

\section{Razporeditev pomnilnika in zanesljiv zagon}

STM32N6570-DK uporablja več ločenih pomnilniških območij. Končne slike ni bilo
mogoče obravnavati kot eno samo datoteko, ki bi jo zapisali na poljuben naslov.
Aplikacija, uteži modela, začasne aktivacije in zaslonska medpomnilnika imajo
različne zahteve glede velikosti, hitrosti in dostopa. Njihova dejanska
razporeditev je povzeta v tabeli \ref{tab:pomnilniska-razporeditev}; podrobna
izkoriščenost posameznih območij je ovrednotena v poglavju
\ref{ch:rezultati}.

\begin{table}[htbp]
  \centering
  \caption{Razporeditev glavnih programskih delov po pomnilniku.}
  \label{tab:pomnilniska-razporeditev}
  \renewcommand{\arraystretch}{1.14}
  \small
  \begin{tabular}{>{\raggedright\arraybackslash}p{0.25\textwidth}>{\raggedright\arraybackslash}p{0.25\textwidth}>{\raggedright\arraybackslash}p{0.35\textwidth}}
    \toprule
    Vsebina & Naslovno območje ali začetek & Namen \\
    \midrule
    Podpisana aplikacija & \texttt{0x70100000} & slika programa v zunanjem bliskovnem pomnilniku \\
    Uteži Neural-ART & \texttt{0x70180000} & ločena binarna datoteka prevedenih uteži \\
    Izvedbena slika aplikacije & \texttt{0x34000400--}\newline\texttt{0x34100000} & izvajanje programa iz notranjega pomnilnika \\
    Začasne aktivacije modela & \texttt{0x34350000--}\newline\texttt{0x343C0000} & namensko notranje območje za Neural-ART \\
    Prvi zaslonski medpomnilnik & \texttt{0x90E80000} & trenutno prikazana slika ali slika v pripravi \\
    Drugi zaslonski medpomnilnik & \texttt{0x90F3B800} & skrita slika za menjavo brez trganja \\
    \bottomrule
  \end{tabular}
\end{table}

Po ponastavitvi se najprej izvede nespremenjen začetni nalagalnik podjetja ST.
Ta v zunanjem bliskovnem pomnilniku poišče podpisano aplikacijo, preveri njeno
glavo, jo naloži v notranji pomnilnik in ji preda izvajanje. Uteži ostanejo v
zunanjem bliskovnem pomnilniku in so dostopne na naslovih, ki jih je določil
prevajalnik modela.

Pri prvem poskusu podpisovanja se program ni zagnal, čeprav je bilo zapisovanje
v pomnilnik uspešno. Vzrok ni bil v modelu ali mikrofonu, ampak v velikosti
predpone podpisane slike. Orodje je privzeto ustvarilo predpono dolžine
\texttt{0x240}, uporabljeni povezovalni opis in začetni nalagalnik pa sta
pričakovala tabelo prekinitvenih vektorjev pri odmiku \texttt{0x400}. Po
nastavitvi poravnave na \texttt{0x400} je podpisana slika vsebovala pravilno
1.024-bajtno predpono in zagon je postal ponovljiv. Ta primer pokaže, da uspešno
prevajanje in zapisovanje še ne zagotavljata zagonsko veljavne slike.

Dodaten problem se je pojavil pri vključitvi zaslona. Dva zaslonska
medpomnilnika sta skupaj prevelika za smiselno namestitev ob aplikaciji in
aktivacijah modela v notranjem pomnilniku, zato sta bila prestavljena v zunanji
pomnilnik HyperRAM. Območji sta izbrani na vrhu prvega 16-mebibajtnega dela
pomnilnika in se ne prekrivata z omrežjem. Vsak medpomnilnik vsebuje
$800 \cdot 480 \cdot 2 = 768.000$ bajtov; skupaj to pomeni 1.536.000 bajtov.

Pri nastavitvi prikazovalnika je bilo treba urediti tudi dovoljenja za dostop
krmilnika zaslona do zunanjega pomnilnika in njegov vir ure. Za uporabljene
dele krmilnika so bila nastavljena varna privilegirana dovoljenja. Četrta
fazno sklenjena zanka je bila nastavljena na \SI{50}{\mega\hertz}, vmesni
delilnik pa je pripravil slikovno uro \SI{25}{\mega\hertz}. Pred popravkom je
krmilnik izbiral ta vir, čeprav v izhodiščni zvočni aplikaciji ni bil vključen,
zato zaslon ni deloval pravilno. Diagnostika je tako potekala po posameznih
plasteh: najprej veljavnost zunanjega pomnilnika, nato pravice dostopa in na
koncu izvor ure.

\section{Časovno filtriranje in zavračanje negotovih vhodov}

Izhod nevronskega modela je šest ocen za trenutni zvočni blok. Največja ocena
se lahko med zaporednimi bloki spremeni tudi pri istem dogodku, zlasti kadar je
zvok tih, kratek ali podoben drugemu razredu. Neposreden prikaz največjega
izhoda bi zato povzročal preskakovanje napisov in veliko kratkih alarmov. Med
model in uporabniški vmesnik smo dodali odločitveni filter z glajenjem,
histerezo in posebnim pravilom za govor.

Za vsak razred najprej izračunamo eksponentno drseče povprečje

\begin{equation}
  s_t = 0{,}65 p_t + 0{,}35 s_{t-1},
  \label{eq:eksponentno-glajenje}
\end{equation}

kjer je $p_t$ trenutna izhodna ocena modela, $s_t$ pa glajena ocena. Pri prvem
aktivnem bloku se glajene ocene inicializirajo neposredno z izhodi modela. Po
izračunu se razredi uredijo po padajoči glajeni oceni; tri najvišje vrednosti
se posredujejo zaslonu in diagnostičnemu izhodu.

Tabela \ref{tab:odlocitveni-pragi} prikazuje končne nastavitve filtra. Za
potrditev novega nevarnostnega razreda mora njegova glajena ocena doseči 0,65.
Ko je razred že potrjen, ostane aktiven do padca pod 0,50. Različna pragova
ustvarita histerezo, ki preprečuje preklapljanje tik ob eni mejni vrednosti.
Če želi drug razred zamenjati že potrjenega, mora poleg vstopnega praga
presegati njegovo oceno za najmanj 0,08.

\begin{table}[htbp]
  \centering
  \caption{Končne nastavitve odločitvenega filtra.}
  \label{tab:odlocitveni-pragi}
  \renewcommand{\arraystretch}{1.16}
  \begin{tabular}{>{\raggedright\arraybackslash}p{0.55\textwidth}>{\raggedright\arraybackslash}p{0.24\textwidth}}
    \toprule
    Pravilo & Vrednost \\
    \midrule
    Delež novega bloka v eksponentnem drsečem povprečju & 0,65 \\
    Prag za potrditev nevarnostnega razreda & 0,65 \\
    Prag za sprostitev potrjenega razreda & 0,50 \\
    Najmanjša prednost za neposredno zamenjavo razreda & 0,08 \\
    Prag zaščitnega razreda govor & 0,27 \\
    Dokazni prag za nevihto & 0,32 \\
    Število zaporednih aktivnih blokov za dokaz nevihte & 2 \\
    Prag za sprostitev potrjene nevihte & 0,25 \\
    Število zaporednih neaktivnih blokov za stanje čakanja & 2 \\
    \bottomrule
  \end{tabular}
\end{table}

Govor uporablja asimetrično pravilo. Če njegova glajena ocena doseže 0,27,
lahko takoj prekliče nevarnostno odločitev, čeprav še ni dosegel običajnega
praga 0,65. Nižji prag je bil določen na razvojnem naboru z namenom
zmanjševanja napačnih nevarnostnih alarmov zaradi človeškega glasu. Potrjen
govor je le informacija: prikaže se kot prepoznan zvok, vendar nikoli ne
zaklene nevarnostnega opozorila.

Nevihta uporablja ločeno dokazno pravilo, ker je bila po prenosu večjega
modela na razvojni plošči pogosto pravilno uvrščena na prvo mesto, vendar z
nižjo oceno od skupnega vstopnega praga. Kandidat za nevihto mora v dveh
zaporednih aktivnih blokih doseči najmanj 0,32, na potrditvenem bloku pa mora
biti tudi najvišje uvrščen razred. Potrjena nevihta se sprosti pod vrednostjo
0,25. Pravilo je bilo določeno na ločenem kalibracijskem naboru s sedmimi
predvajanji nevihte in osmimi negativnimi predvajanji. Nato je bilo zamrznjeno
pred sprejemnim preskusom, zato se njegova nastavitev ne spreminja na podlagi
končnih rezultatov.

Filter razlikuje tri vrste končnega stanja. Stanje čakanja pomeni, da vhodni
nadzor aktivnosti v dveh zaporednih blokih ni zaznal dovolj energije. Stanje
neznano pomeni, da je zvok aktiven, vendar noben nevarnostni razred ne doseže
vstopnega praga in govor ne doseže svojega zaščitnega praga. Tretja možnost je
potrjen razred. Ta razlika je pomembna: neznano ni sedmi razred nevronske
mreže, temveč zavrnitev nezanesljive napovedi, čakanje pa označuje neaktiven
vhod.

Uporabniški vmesnik namenoma loči trenutno napoved od zgodovine dogodkov.
Glavna kartica ob aktivnem zvoku sproti kaže najvišji glajeni izhod, zato se
odzove že na spremembo trenutne napovedi. V zgodovino pa se vpiše potrjen
dogodek; sprememba razreda znotraj iste zvočne seje se sprejme po dveh
zaporednih prikaznih oknih. Tako je glavni prikaz odziven, seznam nedavnih
dogodkov in zaklenjeni alarm pa sta stabilnejša.

\section{Uporabniški vmesnik in diagnostični izhod}

Za končni prototip smo izdelali lasten, neposreden uporabniški vmesnik. Pri
izvedbi ni uporabljen TouchGFX. Procesorsko jedro riše osnovne like, črte in
besedilo neposredno v slikovni medpomnilnik, krmilnik zaslona pa ta pomnilnik
prikazuje na plošči. Ta rešitev ima manj odvisnosti in omogoča natančen nadzor
nad trenutkom menjave slike.

Zaslon ima ločljivost $800 \times 480$ slikovnih točk. Vsaka slikovna točka je
shranjena v 16-bitnem barvnem zapisu RGB565, ki rdeči in modri komponenti
nameni po pet bitov, zeleni pa šest. Uporabljena sta dva slikovna
medpomnilnika. Program celoten nov prikaz nariše v trenutno skriti
medpomnilnik, očisti ustrezno območje podatkovnega predpomnilnika in nato ob
navpičnem premoru zaslona zamenja prikazani naslov. S tem se odpravi trganje in
večina vidnega utripanja, ki se je pojavljalo pri risanju neposredno v trenutno
prikazano sliko.

Glavna kartica prikazuje trenutno zaznani zvok, besedilo o stanju ter odstotno
oceno in vodoravni stolpec. Pod njo so tri najvišje glajene napovedi, na desni
pa tri nedavne potrjene zaznave z oceno in časom. V glavi je večji čas
delovanja naprave. Spodnja vrstica prikazuje zaklenjeno opozorilo oziroma
navodilo za upravljanje. Fizični uporabniški gumb začasno ustavi ali nadaljuje
obdelavo, fizični gumb z oznako TAMP pa potrdi in izbriše zaklenjeno opozorilo.
Oznaka TAMP je tovarniška oznaka gumba za zaznavanje posega; v tej aplikaciji
je gumb dobil novo, varno programsko vlogo potrditve opozorila.

Od odstotka na zaslonu ne smemo sklepati o pravilnosti celotnega sistema.
Prikazana vrednost je glajena izhodna ocena modela za trenutni zvočni blok.
Na primer, prikaz 90 odstotkov pomeni močno notranjo oceno za določen razred,
ne pa dokaza, da je sistem v 90 odstotkih vseh resničnih primerov pravilen.
Pravilnost je zato določena s posebej zasnovanim fizičnim vrednotenjem v
naslednjih poglavjih.

Slika \ref{fig:ui-stanji} prikazuje dve značilni stanji končne aplikacije. V
levem primeru ni dovolj aktivnega zvoka in sistem čaka. V desnem primeru je
potrjena sirena, zato so rezultat, ocena in spodnje opozorilo poudarjeni z
barvo nevarnosti. Slika \ref{fig:prototip-okolje} prikazuje delujoči sistem v
namiznem preskusnem okolju, v katerem so bili predvajani zvočni posnetki.

\begin{figure}[htbp]
  \centering
  \begin{subfigure}[b]{0.48\textwidth}
    \centering
    \includegraphics[width=\textwidth]{20260809_185954.jpg}
    \caption{Stanje čakanja.}
  \end{subfigure}
  \hfill
  \begin{subfigure}[b]{0.48\textwidth}
    \centering
    \includegraphics[width=\textwidth]{20260809_185928.jpg}
    \caption{Potrjena zaznava sirene.}
  \end{subfigure}
  \caption{Končni uporabniški vmesnik med čakanjem in po zaznanem dogodku.}
  \label{fig:ui-stanji}
\end{figure}

\begin{figure}[htbp]
  \centering
  \includegraphics[width=0.78\textwidth]{20260809_185913.jpg}
  \caption{Delujoči prototip v namiznem preskusnem okolju.}
  \label{fig:prototip-okolje}
\end{figure}

Poleg zaslona aplikacija uporablja asinhrono serijsko povezavo s hitrostjo
14.400 bitov na sekundo. Ob spremembi odločitve pošlje kratko strukturirano
sporočilo v predmetni notaciji jezika JavaScript, za vsak obdelani zvočni blok pa še strukturirano
vrstico z oznako \texttt{AED\_CSV}. Vrstica vsebuje zaporedno številko bloka,
stanje vhodne aktivnosti, vsoto spektrograma, končno odločitev in njeno oceno,
tri najvišje razrede z ocenami ter podatek, ali se je odločitev spremenila.

Serijski izhod ni namenjen le razhroščevanju. Orodja na osebnem računalniku so
ga med fizičnimi preskusi zajemala neposredno v datoteke, iz katerih so bili
pozneje izračunani rezultati. S tem se nismo zanašali na ročno prepisovanje
napisov z zaslona. Model, prevedene uteži, podpisana aplikacija, preskusni
manifesti in posamezni zvočni dražljaji imajo zapisane 256-bitne kriptografske
kontrolne vsote. Različico na plošči dodatno določa identifikator spremembe v
repozitoriju. Tako je mogoče meritev povezati z natančno programsko in
modelsko različico brez sklepanja iz imena datoteke.

Implementacija tako tvori neprekinjeno verigo: vgrajeni mikrofon zajame zvok,
strojni filter in neposredni prenos ga pretvorita v šestnajstbitne vzorce,
predobdelava izračuna kvantizirani logaritemski melov spektrogram, Neural-ART
izvede prevedeni model, časovni filter pripravi stabilnejšo odločitev, rezultat
pa se hkrati prikaže uporabniku in zabeleži za ponovljivo vrednotenje.
